% !TEX root = ../nc-cubic.tex
\section{中心一次変数が一つの混合ブロックと有限 norm}
\label{cmo:sec:points}

分離表示の定理\ref{csc:thm:comparison}では、中心変数は各楕円因子の
二次関係に現れない。本節では、この仮定を点関手の計算から取り除く。
結果は中心変数を一つ共有する表示についての定理であり、
一般の混合表示の点関手を計算した後、$S(1,-1,-1)$ の定数型変形については
有限中心性と収縮族の相対 smoothness を証明し、余有限可換 Ku 実現を得る。

\subsection{有限な定常環と point scheme}

基礎体を $\C$、$R=\C[u]$、$k\geq2$ とする。
各 $i$ について、次数一生成元 $x_i,y_i,z_i,u$ を持ち、
$u$ が中心である二次代数 $\mathcal B_i$ を取る。
$\mathcal B_i/(u)=S_i$ の双線型対応は、滑らかな平面三次 $E_i$ 上の
固定点なしの自己同型 $\sigma_i$ のグラフであると仮定する。
さらに可換化は
\begin{equation}\label{cmo:eq:abelian}
 (\mathcal B_i)_{\mathrm{ab}}
   =R[x_i,y_i,z_i]/(q_{i1},q_{i2},q_{i3}),\qquad
 q_{ia}\in R[x_i,y_i,z_i]_2
\end{equation}
であり、$q_{i1}|_{u=0},q_{i2}|_{u=0},q_{i3}|_{u=0}$ は
$\C[x_i,y_i,z_i]$ の正則列とする。
例えば、それらが滑らかな三次式の三つの偏微分ならばこの条件を満たす。
active 変数と $u$ の混合項は許す。

異なるブロックの元を可換とする $R$ 上の tensor 積
\[
 B=\mathcal B_1\otimes_R\cdots\otimes_R\mathcal B_k
\]
を考える。正則中心三次 $w\in Z(B)_3$ は、$u=0$ において
$\sum_i a_i\Omega_i$ に移ると仮定する。
ここで各 $\Omega_i$ は $S_i$ の全点加群を消す中心三次である。
$c\in\C$ に対し $A_c=B/(w+cu^3)$ と置く。
次の定理の点関手の主張自体には、$A_c$ の滑らかさを仮定しない。

\begin{theorem}[一中心変数での定常点の有限判定]
\label{cmo:thm:point-norm}
上の仮定の下で
\[
 T=R[x_1,y_1,z_1,\ldots,x_k,y_k,z_k]/(q_{ia})
\]
は $R$ 上階数 $L=8^k$ の有限自由加群である。
$Q=T/(u-1)$、$b=w_{\mathrm{ab}}|_{u=1}\in Q$ とすると、任意の試験環に対して
\begin{equation}\label{cmo:eq:point-scheme}
 \operatorname{Pt}(B)=\left(\coprod_i E_i\right)\amalg\operatorname{Spec}Q,
 \qquad
 \operatorname{Pt}(A_c)=\left(\coprod_i E_i\right)
                  \amalg\operatorname{Spec}(Q/(b+c)).
\end{equation}
shift は $E_i$ 上で $\sigma_i$、定常成分上で恒等である。

乗法 $m_b:Q\to Q$ の norm 多項式
\begin{equation}\label{cmo:eq:norm}
 N(c)=\det(c\operatorname{id}_Q+m_b)\in\C[c]
\end{equation}
は次数 $8^k$ の monic 多項式であり、
\[
 \operatorname{Pt}(A_c)=\coprod_iE_i
 \quad\Longleftrightarrow\quad N(c)\ne0.
\]
従って定常点が現れるパラメータは有限個で、
その零次 Fitting イデアルは $(N(c))$ である。
これは閉点集合だけでなく、$c$ に関する族の有限表示としての主張である。
\end{theorem}

\begin{proof}
$u,q_{11},\ldots,q_{k3}$ の同時零点は affine 空間の原点だけである。
これらは $3k+1$ 変数多項式環の斉次 system of parameters なので正則列である。
従って $T$ は Hilbert 級数
\[
 H_T(t)=\frac{(1+t)^{3k}}{1-t}
\]
を持つ。$T/uT$ の斉次基底を $T$ へ持ち上げると、
次数についての帰納法によりそれらは $R$-加群として $T$ を生成する。
基底の次数を保つ自由 $R$-加群からの全射の両辺の Hilbert 級数は同じなので、
その核は零である。これで階数 $8^k$ の自由性を得る。
また $C=\operatorname{Proj}T$ は $u=0$ と交わらず、
$u\ne0$ の chart で $\operatorname{Spec}Q$ と同一視される。

双線型関係を満たす二つの点を $p,p'$ と書く。
$u(p)\ne0$ なら、中心交換関係
$u(p)x'_i=x_i u(p')$ などと $p'\ne0$ から $u(p')\ne0$ であり、
両方を一に規格化すると全座標が等しい。
残る関係は式\eqref{cmo:eq:abelian}の $3k$ 本の二次式 $q_{ia}=0$ である。
よってこの chart の対応は $\Delta_C$ である。

$u(p)=0$ の幾何点では $u(p')=0$ となる。
二つ以上のブロックが非零ならブロック間交換関係が各ブロックの二点を比例させ、
$\sigma_i$ の固定点を生じるので不可能である。
従ってその幾何点は純ブロックの $\Gamma_{\sigma_i}$ にある。
そこでは固定点なしという仮定により、同一ブロックの二点の
二行二列小行列式の少なくとも一つが単元となる。
中心交換関係の対応する二式は $u,u'$ を環の中で消し、
他ブロックとの交換関係も同じ単元小行列式によって全ての他座標を消す。
残る対応は $\Gamma_{\sigma_i}$ そのものである。
従って横断方向の冪零構造もなく、
\[
 \Gamma_B=\left(\coprod_i\Gamma_{\sigma_i}\right)\amalg\Delta_C
\]
をスキームとして得る。この局所環上の消去は任意の試験環でも成立する。

対応が自己同型のグラフなので、無限点列は初期点から一意に定まる。
純ブロックでは $u=0$、かつ $\Omega_i$ は普遍点加群を消す。
定常列では三次関係の三重線型化はその可換化の評価であり、
$u=1$ で $b+c$ となる。これで式\eqref{cmo:eq:point-scheme}を得る。

普遍パラメータ $c$ を使うと、$Q[c]$ 上の乗法 $m_b+c$ は単射である。
実際、非零の多項式ベクトルの最高次係数は $c\operatorname{id}$ により残る。
従って
\[
 0\longrightarrow Q[c]\xrightarrow{m_b+c}Q[c]
   \longrightarrow Q[c]/(b+c)\longrightarrow0
\]
は階数 $L$ の正方表示を与え、その零次 Fitting イデアルは行列式である。
閉点 $c$ で商が零となることは $b+c$ が $Q$ の単元であること、
従って $N(c)\ne0$ と同値である。
\end{proof}

\subsection{滑らかな Poisson 標的での定常零点の消滅}

同じ変数分解の下で
\[
 H=\sum_{i=1}^kF_i(x_i,y_i,z_i,u)+G(u),\qquad
 X=V(H)\subset\mathbb P^{3k}
\]
とする。各 $F_i$ は斉次三次であり、$g_i=F_i|_{u=0}$ は滑らかな平面三次を定め、
$X$ も滑らかとする。
$P_i=\partial_{x_i}\wedge\partial_{y_i}\wedge\partial_{z_i}$、
$d_i\in\C^\times$ を取り、
\[
 \beta=\sum_i d_i\iota_{dF_i}P_i
\]
の射影化を $\pi$ と書く。共通の $u$ は各項の Casimir であり、
異なるブロックの Schouten 括弧も零なので $\beta$ は Poisson で、$H$ を保つ。

\begin{proposition}[一中心変数の Poisson 零スキーム]
\label{cmo:prop:poisson-zero}
上の仮定の下で、非被約構造を含めて
\begin{equation}\label{cmo:eq:poisson-zero}
 Z_X(\pi)=\coprod_iV(g_i)
\end{equation}
である。右辺はそれぞれの純ブロック平面に埋め込む。
\end{proposition}
\begin{proof}
Euler vector field を $\xi$ とする。
命題\ref{pps:prop:zero-scheme}と同じく零イデアルは $\xi\wedge\beta$ の係数である。
$u\ne0$ では $u\,d_i(F_i)_{x_i}$ などの係数から、
全ての active 偏微分が零となる。
$H=0$ と Euler 恒等式から $uH_u=0$ も従うので、
もしこの chart に零点があれば $X$ の特異点となる。
従って滑らかさの仮定によりこの chart の零スキームは空である。

$u=0$ の幾何点で $s$ 個のブロックが非零なら、$g_i$ の滑らかさから
$\beta$ の階数は $2s$ である。
射影化が零なら $\beta\in\xi\wedge V$ なので階数は高々二であり、$s=1$ である。
純ブロック $i$ では残る係数は $3d_i g_i$ である。
さらにその点の近傍では $(F_i)_{x_i},(F_i)_{y_i},(F_i)_{z_i}$ の一つが単元なので、
混合外積の係数は $u$ と全ての他ブロック座標を環の中で消す。
残る式は $g_i=0$ となり、式\eqref{cmo:eq:poisson-zero}を得る。
\end{proof}

\subsection{六十四個の相異なる例外値を持つ表示}

次の例では、定常成分が現れる全パラメータを有理数で記述できる。
ここで構成する中心商の各例外値での coherent smoothness は未確認であり、
第二予想の反例としては扱わない。

\begin{example}[混合二ブロックの完全な norm 分解]
\label{cmo:ex:sixty-four}
$\alpha=(\alpha_1,\alpha_2,\alpha_3)$ に対し、$u$を中心とする代数
\[
 \mathcal B_\alpha=\C\langle x,y,z,u\rangle/
 \left([y,z]-x^2+\alpha_1u^2,
       [z,x]-y^2+\alpha_2u^2,
       [x,y]-z^2+\alpha_3u^2,
       [u,x],[u,y],[u,z]\right)
\]
を考える。次の二組を用いる。
\[
 \alpha=(1,4,9),\qquad \alpha'=(16,25,49).
\]
各代数で、$\sym(xyz)$を六語の平均として
\begin{equation}\label{cmo:eq:mixed-center}
 \Omega_\alpha
 =\frac{x^3+y^3+z^3}{3}-\frac23\sym(xyz)
       -u^2(\alpha_1x+\alpha_2y+\alpha_3z)
\end{equation}
は中心元である。
$B=\mathcal B_\alpha\otimes_{\C[u]}\mathcal B_{\alpha'}$、
$w_c=\Omega_\alpha+100\Omega_{\alpha'}+cu^3$と置く。
定理\ref{cmo:thm:point-norm}の$Q$は$\C^{64}$と同型であり、
そのnormは
\begin{equation}\label{cmo:eq:explicit-norm}
 \begin{gathered}
 N(c)=\prod_{r\in\mathcal R}\prod_{s\in\mathcal S}(c+r+100s),\\
 \mathcal R=\{\pm28,\ \pm16,\ \pm56/3,\ \pm28/3\},\qquad
 \mathcal S=\{\pm448,\ \pm196,\ \pm176,\ \pm284/3\}
 \end{gathered}
\end{equation}
となる。六十四個の根は全て相異なる。
各根では点スキームにちょうど一つの被約な定常点が加わり、
根以外では二つの楕円曲線だけとなる。特に$c=0$では定常成分は空である。
\end{example}

\begin{proof}
$u=0$の代数はSklyanin parameters $(1:-1:-1)$を持ち、
点曲線は$E=V(x^3+y^3+z^3+xyz)$である。
$1+(1/3)^3\ne0$より滑らかで、
対角関係$x^2=y^2=z^2=0$は射影共通零点を持たない。
命題\ref{esi:prop:factor}と同じ行列の核の議論で、
双線型対応は固定点なしの楕円グラフである。
このSklyanin代数は\cite[Theorem~1.1, Corollary~1.3]{GRSklyaninTV}によりKoszulである。

$u=1$の関係は定数項だけを加えたPBW変形である。
実際、$r_x=yz-zy-x^2$とその巡回置換を用いると、
唯一の三次Koszul重なりは
\[
 x\otimes r_x+y\otimes r_y+z\otimes r_z
   =r_x\otimes x+r_y\otimes y+r_z\otimes z
\]
で生成される。
Braverman--Gaitsgoryの判定
\cite[Main Theorem~0.5, Lemma~3.3]{CMOBravermanGaitsgory}における
一次補正は零、定数補正は$r_a\mapsto-\alpha_a$である。
残る条件はこの重なり上の
$\beta\otimes1-1\otimes\beta=0$であり、定数の可換性から成立する。
従って表示した斉次代数はこのPBW変形のRees代数で、$\C[u]$上自由である。

式\eqref{cmo:eq:mixed-center}の中心性は、指定した二組の係数について
有理数上の有限な線型代数で直接検証する。
次数$d$の全語を基底に、二次関係を長さ$d-2$の全左右語で挟んだ行列を作る。
$d=3,4$の商次元は$20,35$であり、$\Omega_\alpha$の次数三剰余は非零、
四生成元との全交換子の次数四剰余は零となる。
再現コードは\texttt{verify\_point\_poisson\_obstruction.py}の
\texttt{verify\_mixed\_central\_norm}にあり、全行列を関係式から再生成して厳密消去する。
中心性は生成元とのこの有限個の交換子で尽くされ、
全次数PBW性は前段の判定で別に示している。
$u=0$では、$S=x^3+y^3+z^3$、
$C_+=xyz+yzx+zxy$、$C_-=yxz+xzy+zyx$と書くと
\[
 C_+-C_--S=0,\qquad 3c_3=2C_--2S,\qquad
 \Omega_0=\frac{3S-C_+-C_-}{9}=-\frac13c_3.
\]
ここで標準中心元$c_3$の式は
\cite[Section~2, p.~3]{GRSklyaninDLL}に$(a,b,c)=(1,-1,-1)$を代入したものである。
従って$\Omega_0$は全点加群を消す。
二因子の中心三次和は$B/(u)$上正則であり、
$u$の正則性と次数についての帰納法から、$w_c$も全ての$c$で正則である。

可換化して$u=1$とすると$x^2=\alpha_1,y^2=\alpha_2,z^2=\alpha_3$である。
各因子は八つの相異なる有理点に分解し、
その上の中心元の値は
\[
 b_\alpha=-\frac23(\alpha_1x+\alpha_2y+\alpha_3z+xyz)
\]
となる。$(x,y,z)=(\pm1,\pm2,\pm3)$での全評価値が$\mathcal R$、
$(\pm4,\pm5,\pm7)$での全評価値が$\mathcal S$である。
従って$Q=\C^{64}$上の乗法は対角となり、式\eqref{cmo:eq:explicit-norm}を得る。
$\mathcal S$の異なる二元の差の絶対値は少なくとも$20$、
$\mathcal R$の直径は$56$なので、$100\mathcal S$を中心とする八つの組は交わらない。
各組の八値も相異なるため、全ての根は単純である。
また$100\min_{s\in\mathcal S}|s|>\max_{r\in\mathcal R}|r|$より$N(0)\ne0$である。
点関手の結論は定理\ref{cmo:thm:point-norm}から従う。
\end{proof}

\begin{remark}[二つの予想との接続範囲]
\label{cmo:rem:scope}
$k=2,m=1$ は $N=7,n=5$ に当たり、定理\ref{cmo:thm:point-norm}の
有限定常環の長さは $64$ である。
一般の中心係数では点スキームは二つの楕円曲線だけとなり、
命題\ref{cmo:prop:poisson-zero}の零スキームと同じ形を持つ。
しかし、一般の混合表示について、その楕円曲線を active 制限に持つ
\emph{同じ}可換 Ku 標的を構成したことにはならない。
分離表示では定理\ref{csc:thm:comparison}がこの不足を埋める。

例外的な中心係数で定常点が生じても、それだけで第二予想への新しい反例とはしない。
その係数での非可換三次の滑らかさと、別の Poisson support を含む全標的の排除が必要である。
有限 norm による明示的な例外集合は、これらの次の検証対象を与える。
\end{remark}

\subsection{例外的な定常点の自己 Ext と中心方程式の高次作用}
\label{cmo:subsec:character-ext}

定常点の出現と商の非滑らかさは同じ条件ではない。
本節の明示族では、例外値の定常 character の自己 Ext を全次数で計算できる。
結論は自己 Ext の次数付きベクトル空間についてであり、
その Yoneda 積、character の射影次元、商の smoothness は区別する。

高次 homotopy から超曲面商上の自由分解を作る一般構成は
Shamash--Eisenbud 理論に属する
\cite[Theorem~7.2]{NFMEisenbudCI}。
この可換環上の構成と、本稿の正則中心元に対する非可換版を区別し、
後者の完全性は補題\ref{mbnp:lem:shamash}の濾過による証明を用いる。
また非斉次三次 potential の Calabi--Yau 性、PBW 性と中心元は
Etingof--Ginzburg の既知の理論
\cite[Theorem~3.3.2, Sections~3.4,~8.3]{NFMEG}に位置付けられる。
本節では指定中心元の普遍な多項式証明書を作り、有限消去から
式\eqref{cmo:eq:kappa}を導いて六十四個の自己 Ext を全次数で決定する。
一般の分解構成や周期性自体の新規性は主張せず、
この具体式と自己 Ext 計算の包括的な新規性確認も未完了である。

$p=(a,b,c)\in(\C^\times)^3$ とし、
$\alpha=(a^2,b^2,c^2)$、
$D_p=\mathcal B_\alpha/(u-1)$ と置く。
$M_p=\C$ は $(x,y,z)$ が $(a,b,c)$ として作用する左 $D_p$-加群である。
式\eqref{cmo:eq:mixed-center}の $u=1$ への特殊化を $\Omega_p$ と書き、
\begin{equation}\label{cmo:eq:character-value}
 \lambda_p=-\frac23(a^3+b^3+c^3+abc),\qquad
 f_p=\Omega_p-\lambda_p,\qquad R_p=D_p/(f_p)
\end{equation}
とする。$f_p$ は $M_p$ を消す正則中心元である。
ここではこの中心元の正規化を固定する。

\begin{proposition}[character における中心三次の二次作用]
\label{cmo:prop:character-ext}
次の有理式を置く。
\begin{equation}\label{cmo:eq:kappa}
 \kappa(a,b,c)=
 \frac{2\bigl((ab)^3+(bc)^3+(ca)^3+9a^2b^2c^2\bigr)}{9abc}.
\end{equation}
自己 Ext を計算する cochain 複体として
\begin{equation}\label{cmo:eq:character-complex}
 \operatorname{RHom}_{R_p}(M_p,M_p)
  \simeq
  \bigl(\C[t]\oplus\C[t]\eta,\quad
        |t|=2,\ |\eta|=3,\ d(t)=0,\ d(\eta)=\kappa(p)t^2\bigr)
\end{equation}
がある。この表示の右辺の乗法を Yoneda 積と同定する主張は含めない。
従って $\kappa(p)\ne0$ なら
\[
 \operatorname{Ext}_{R_p}^j(M_p,M_p)=
 \begin{cases}\C,&j=0,2,\\0,&\text{その他},\end{cases}
\]
である。$\kappa(p)=0$ なら、自己 Ext は次数 $2q$ と $2q+3$
（$q\geq0$）に各一次元あり、上に有界でない。
中心方程式を $s f_p$、$s\in\C^\times$ に変更したとき、
同じ Shamash 変数の規約では式\eqref{cmo:eq:character-complex}の係数は
$s^2\kappa(p)$ となる。
\end{proposition}

\begin{proof}
$X=x-a,Y=y-b,Z=z-c$ と置く。関係式は
\[
 [Y,Z]=X^2+2aX,\qquad [Z,X]=Y^2+2bY,\qquad
 [X,Y]=Z^2+2cZ
\]
となる。左自由加群の係数を行として書き、写像の行列を右から掛ける規約で、
\begin{equation}\label{cmo:eq:character-resolution}
 0\longrightarrow D_p\xrightarrow{d_3}D_p^3
   \xrightarrow{d_2}D_p^3\xrightarrow{d_1}D_p
   \longrightarrow M_p\longrightarrow0,
 \quad
 \begin{gathered}
 d_1=(X,Y,Z)^{\mathsf t},\qquad d_3=(X,Y,Z),\\
 d_2=\begin{pmatrix}
  -X-2a&-Z&Y\\ Z&-Y-2b&-X\\ -Y&X&-Z-2c
 \end{pmatrix}
 \end{gathered}
\end{equation}
を得る。$d_2d_1=0$ と $d_3d_2=0$ は表示関係そのものである。
第 $i$ 項を $i$ だけずらした通常の次数濾過を取ると、associated graded は
$S(1,-1,-1)$ の augmentation の Koszul 分解である。
例\ref{cmo:ex:sixty-four}の PBW 論証と Koszul 性により、
この有界下濾過複体も完全である。
$\Omega_p$ の最高次項は正則中心元 $\Omega_0$ なので、
同じ濾過から $f_p$ の正則性を得る。
中心性は、以下の普遍な有限証明書でも同時に確認する。

$P_\bullet$ を式\eqref{cmo:eq:character-resolution}の自由分解とする。
$f_p$ の乗法に対する higher homotopies を
\[
 h_{1,i}:P_i\longrightarrow P_{i+1}\quad(0\leq i\leq2),
 \qquad h_2:P_0\longrightarrow P_3
\]
と書く。$\mathbb Q[a,b,c]$ 上で、係数がそれぞれ全次数二、三の行列を取り、
\begin{equation}\label{cmo:eq:character-homotopy}
 dh_1+h_1d=f_p\operatorname{id},\qquad
 dh_2+h_2d+h_1^2=0
\end{equation}
を満たすようにできる。長さ三の分解なので、これより高い恒等式は
写像先の次数の理由で零になる。
ここで使う具体的な多項式行列の全係数と、その有限証明書の検算を
\texttt{computations/verify\_mixed\_character\_ext.py} に収めた。
三つの中心パラメータも次数一とし、全次数 $2,3,4$ の全語に対し、
各関係の全左右語倍と全パラメータ倍を列挙する。
その有理係数関係行列で、$d^2=0$、$[f_p,X]=[f_p,Y]=[f_p,Z]=0$、
式\eqref{cmo:eq:character-homotopy}を全て零へ還元する。
この検算は有限個のパラメータの代入や補間ではなく、
$\mathbb Q[a,b,c]$ 上の多項式恒等式である。
全次数での分解の完全性は前段の PBW/Koszul 論証で別に保証した。

この証明書の character での値だけを記すと、
\begin{equation}\label{cmo:eq:character-homotopy-values}
 \begin{gathered}
 d_2(p)=\operatorname{diag}(-2a,-2b,-2c),\qquad h_{1,1}(p)=0,\\
 v=h_{1,0}(p)=-\frac23(bc,ac,ab+c^2),\\
 z=h_{1,2}(p)=\frac23(a^2-bc,-ac-b^2,-ab)^{\mathsf t},\qquad
 h_2(p)=\frac{16}{9}abc.
 \end{gathered}
\end{equation}
補題\ref{mbnp:lem:shamash}の分解構成は、内部次数を忘れた正則中心元と
任意の消される加群にも同じ証明で適用できる。
すなわち $\bigoplus_{q\geq0}(R_p\otimes_{D_p}P_{n-2q})$ に
$d,h_1,h_2$ を配置する。添字 $q$ による濾過の第一頁は
二項中心分解から $M_p$ を二つ持ち、$dh_1+h_1d=f_p$ による連結射が
隣接する項を恒等で相殺する。各全次数の濾過は有限なので、
これは $M_p$ の自由 $R_p$-分解となる。

この分解へ $\operatorname{Hom}_{R_p}(-,M_p)$ を施すと、
\[
 \left(\operatorname{Hom}_{D_p}(P_\bullet,M_p)[t],
       d^*+t h_1^*+t^2h_2^*\right)
\]
を得る。式\eqref{cmo:eq:character-homotopy-values}の可逆な $d_2(p)$ を
$\C[t]$ 上で消去する。$h_{1,1}(p)=0$ なのでこの消去は有限の行列操作であり、
形式冪級数への完備化を必要としない。
残る次数零と三の二つの基底間の係数は
\[
 \left(h_2(p)-v\,d_2(p)^{-1}z\right)t^2
   =\kappa(p)t^2
\]
である。最後の等式は式\eqref{cmo:eq:character-homotopy-values}の代入による
有理式恒等式であり、同じコードで記号的にも確認する。
これが式\eqref{cmo:eq:character-complex}を与える。
$\kappa\ne0$ なら残る cohomology は $\C[t]/(t^2)$、
$\kappa=0$ なら微分は零である。
また中心元を $s$ 倍すると式\eqref{cmo:eq:character-homotopy}では
$h_1$ を $s$ 倍、$h_2$ を $s^2$ 倍するので、係数は $s^2\kappa$ となる。
\end{proof}

\begin{corollary}[六十四例外点の自己 Ext は有界である]
\label{cmo:cor:sixty-four-ext}
例\ref{cmo:ex:sixty-four}の二ブロックについて、
$p=(\pm1,\pm2,\pm3)$、$p'=(\pm4,\pm5,\pm7)$ とし、
$c_0=-\lambda_p-100\lambda_{p'}$ と置く。
中心 chart の代数
\[
 D_{c_0}=(A_{c_0}[u^{-1}])_0
    \simeq (D_p\otimes D_{p'})/(f_p+100f_{p'})
\]
の character $M=M_p\otimes M_{p'}$ に対して、
\begin{equation}\label{cmo:eq:sixty-four-ext}
 \operatorname{Ext}_{D_{c_0}}^j(M,M)=
 \begin{cases}\C,&j=0,2,3,5,\\0,&\text{その他}\end{cases}
\end{equation}
である。従って全六十四例外値について、この定常 character の自己 Ext は
商の非滑らかさを示す「無限次数の自己 Ext」という障害を与えない。
\end{corollary}

\begin{proof}
二因子の式\eqref{cmo:eq:character-resolution}を tensor し、
homological Koszul 符号付きに higher homotopies を足す。
異なる因子の奇写像は反交換するため、
中心 $f_p+100f_{p'}$ の higher homotopies は、第二因子の
$h_1$ を $100$ 倍、$h_2$ を $10000$ 倍して得られる。
先の $\C[t]$ 上の有限消去を各因子に施すと、
自己 Ext を計算する複体は、次数付きベクトル空間として
\begin{equation}\label{cmo:eq:two-character-complex}
 \left(\C[t]\otimes\bigwedge(\eta_1,\eta_2),\quad
 |t|=2,\ |\eta_i|=3,\quad
 d(\eta_1)=\kappa(p)t^2,\quad
 d(\eta_2)=10000\kappa(p')t^2\right)
\end{equation}
となる。外積と Leibniz 則はこの計算複体の記法であり、
自己 Ext の Yoneda 積の同定は必要ない。

式\eqref{cmo:eq:kappa}の分子で、全ての符号の選び方に対して
\[
 9a^2b^2c^2=324>8+216+27=251
\]
が成立する。第二因子でも
\[
 9a^2b^2c^2=176400>8000+42875+21952=72827
\]
なので、両方の $\kappa$ は零でない。
$\eta=\eta_1/\kappa(p)$ と
$\zeta=\eta_2-10000\kappa(p')\eta_1/\kappa(p)$ を取ると
$d\eta=t^2,d\zeta=0$ である。従って全次数の cohomology は
$\C[t]/(t^2)\otimes\bigwedge(\zeta)$ となり、
式\eqref{cmo:eq:sixty-four-ext}を得る。
全十六因子 character と六十四組の係数・有限行列の検算も
上記再現コードに含めた。全次数の消滅は、この基底変換による証明であり、
有限次数までの計算からの外挿ではない。
\end{proof}

\begin{remark}[自己 Ext 判定の射程]
\label{cmo:rem:self-ext-scope}
可換化の norm $N(c)$ は定常 character の存在を検出する。
有限自己 Ext から有限射影次元を結論する一般論は使えないが、
以下では有限 ambient 分解と二周期性を追加して射影次元まで決定する。
その結論も、全加群の有限大域次元や全 coherent 圏の smoothness の判定とは異なる。
また計算複体に書いた外積を Yoneda 積と同一視する主張は含めない。
\end{remark}

\subsection{周期性を使った character の射影次元判定}

有限自己 Ext だけから有限射影次元を結論することは一般にはできない。
ここでは有限 ambient 分解から得られる二周期性と、中心商の Rees 同型を
追加してこの推論を正当化する。以下の一般論自体の新規性は主張しない。
Shamash 型分解は補題\ref{mbnp:lem:shamash}の証明を用いる。
非可換な正則 normal 商の周期分解については Cassidy--Conner--Kirkman--Moore
\cite[Theorem~4.2]{CMOPeriodicMF} も関連する。
ただし同論文の主定理は連結次数付きの設定なので、以下の affine 版は直接証明する。

\begin{lemma}[二周期性による有限射影次元の検出]
\label{cmo:lem:periodic-pd}
$R$ を環、$M$ を有限生成左 $R$-加群とする。
$M$ が有限生成射影加群による分解を持ち、その十分高い syzygy が
その分解の中で実際に二周期であると仮定する。
さらに $\operatorname{Ext}_R^i(M,R)=0$ が十分大きい $i$ で成立するとする。
このとき
\[
 \operatorname{pd}_R M<\infty
 \quad\Longleftrightarrow\quad
 \operatorname{Ext}_R^i(M,M)=0\quad(i\gg0).
\]
また $d\ge1$ とし、$\operatorname{Ext}_R^i(M,R)=0$ が $i\ge d$ で成立するなら、
有限の場合には $\operatorname{pd}_R M\le d-1$ である。
\end{lemma}
\begin{proof}
順方向は定義から従う。逆方向を示す。
$K=\Omega_R^n M$ を二周期部分に入った syzygy とする。
$M$ の分解の第 $i$ 項を $P_i$ と書く。
各 $P_i$ は有限生成射影なので、$\operatorname{Ext}_R^{\gg0}(M,P_i)=0$ である。
十分大きい $r$ に対して、第一変数、次いで第二変数の dimension shifting から
\[
 \operatorname{Ext}_R^{2r}(K,K)
 \simeq\operatorname{Ext}_R^{2r+n}(M,K)
 \simeq\operatorname{Ext}_R^{2r}(M,M)=0
\]
を得る。第二同型は有限個の完全列
$0\to\Omega_R^{i+1}M\to P_i\to\Omega_R^iM\to0$
に $\operatorname{Hom}_R(M,-)$ を適用して得られる。

$K$ の二周期分解における同型
$\theta:\Omega_R^{2r}K\xrightarrow{\sim}K$ を取る。
標準の表示
\[
 \operatorname{Ext}_R^{2r}(K,K)
 =\operatorname{coker}\!\left(
  \operatorname{Hom}_R(P_{n+2r-1},K)
  \longrightarrow\operatorname{Hom}_R(\Omega_R^{2r}K,K)\right)
\]
において $[\theta]=0$ なので、$\theta$ は $P_{n+2r-1}\to K$ へ延長する。
この延長と $\theta^{-1}$ の合成は
$\Omega_R^{2r}K\hookrightarrow P_{n+2r-1}$ の retraction となる。
従って $K\simeq\Omega_R^{2r}K$ は射影で、$M$ の射影次元は有限である。

最後に $m=\operatorname{pd}_R M<\infty$ とする。
もし $m>0$ かつ $\operatorname{Ext}_R^m(M,R)=0$ なら、長さ $m$ の
有限生成射影分解を双対化した最後の射
$P_{m-1}^{\vee}\to P_m^{\vee}$ は全射である。
$P_m^{\vee}$ は有限生成右射影加群だからこの全射は分裂する。
有限生成射影加群の bidual 同型を使うと
$P_m\to P_{m-1}$ が分裂単射となり、分解を短縮できて矛盾する。
従って $m>0$ なら $\operatorname{Ext}_R^m(M,R)\ne0$。
これより $m\le d-1$ を得る。
\end{proof}

\begin{corollary}[有限 ambient 分解を持つ中心超曲面加群]
\label{cmo:cor:hypersurface-pd}
$Q$ を環、$f\in Z(Q)$ を左右正則元、$R=Q/(f)$ とする。
$M\ne0$ を左 $R$-加群とし、$Q$ 上で長さ $d\ge1$ の
有限生成射影分解を持つと仮定する。
このとき前補題の同値と有限の場合の上限 $\operatorname{pd}_R M\le d-1$ が成立する。
左右を交換した主張も成立する。
\end{corollary}
\begin{proof}
中心元の二項分解から
\[
 \operatorname{RHom}_Q(R,Q)\simeq R[-1],\qquad
 \operatorname{RHom}_Q(M,Q)\simeq
 \operatorname{RHom}_R(M,R)[-1]
\]
を得る。すなわち Rees 同型
\[
 \operatorname{Ext}_R^i(M,R)\simeq
 \operatorname{Ext}_Q^{i+1}(M,Q)=0\quad(i\ge d)
\]
がある。ここで正則元の中央性により normalizing twist は現れない。

$Q$ 上の有限分解 $P_\bullet$ に $f$ の higher homotopies $s_a$ を選ぶ。
$s_a:P_j\to P_{j+2a-1}$ なので非零な $s_a$ は有限個である。
補題\ref{mbnp:lem:shamash}と同じ濾過証明により
\[
 C_n=\bigoplus_{q\ge0}(R\otimes_QP_{n-2q})
\]
は $M$ の有限生成射影分解になる。
十分大きい偶数 $n$ では $C_n=\bigoplus_{j\text{ even}}R\otimes_QP_j$、
十分大きい奇数 $n$ では奇数 $j$ の直和となり、微分も parity のみに依存する。
従ってこれは実際の二周期分解である。
ここでは内部次数を忘れた affine 加群を扱うので、内部次数 shift を
同型と取り違えているのではない。前補題が適用できる。
\end{proof}

\begin{proposition}[混合 character の射影次元の完全判定]
\label{cmo:prop:character-pd}
$p_i=(a_i,b_i,c_i)\in(\C^\times)^3$、$\beta_i\in\C^\times$ とし、
命題\ref{cmo:prop:character-ext}の記法で
\[
 Q=\bigotimes_{i=1}^kD_{p_i},\quad
 R=Q\Big/\Big(\sum_{i=1}^k\beta_if_{p_i}\Big),\quad
 M=\bigotimes_{i=1}^kM_{p_i}
\]
と置く。中心和は正則とする。この仮定は本節の混合族では PBW 濾過から成立する。
このとき
\[
 \operatorname{pd}_R M=
 \begin{cases}
  3k-1,&\text{少なくとも一つの }\kappa(p_i)\ne0,\\
  \infty,&\text{全ての }\kappa(p_i)=0.
 \end{cases}
\]
特に一因子の $\kappa\ne0$ では射影次元は $2$、
例\ref{cmo:ex:sixty-four}の六十四個の affine character では全て $5$ である。
\end{proposition}
\begin{proof}
各因子の三段自由分解を tensor すれば、$M$ は $Q$ 上で長さ $3k$ の
有限自由分解を持つ。二因子の証明と同じ有限消去により、自己 Ext の計算複体は
\[
 \C[t]\otimes\bigwedge(\eta_1,\ldots,\eta_k),\quad
 |t|=2,\quad|\eta_i|=3,\quad d\eta_i=\beta_i^2\kappa(p_i)t^2
\]
となる。全ての $\kappa(p_i)$ が零なら微分は零で、自己 Ext は無限次数にある。
従って射影次元は無限である。
少なくとも一つが非零なら、線型基底変換で
$d\eta=t^2,d\zeta_1=\cdots=d\zeta_{k-1}=0$ となる。
cohomology は次数付きベクトル空間として
\[
 \C[t]/(t^2)\otimes\bigwedge(\zeta_1,\ldots,\zeta_{k-1})
\]
であり、上に有界で最高非零次数は $3k-1$ である。
前系より射影次元は有限で $3k-1$ 以下、最高自己 Ext の非消滅より以上でもある。
\end{proof}

\subsection{中心 chart の character 自己 Ext と qgr の比較}

以下では右加群の規約を使う。先の左 character の計算を移すときは、
有限次元 $M$ に対する自然なベクトル空間同型
\[
 \operatorname{Ext}_{R\text{-left}}^i(M,M)
 \simeq \operatorname{Ext}_{R\text{-right}}^i(M^*,M^*)
\]
を使える。実際、左射影分解 $P_\bullet\to M$ の体上の双対は
右単射分解 $M^*\to P_\bullet^*$ になる。
$P_i$ が平坦であることと tensor--Hom 随伴から $P_i^*$ の単射性が従い、
$\operatorname{Hom}_{R\text{-right}}(M^*,P_i^*)
\simeq\operatorname{Hom}_{R\text{-left}}(P_i,M^{**})$
で比較できる。一文字 character の場合は $M^*$ の座標も同じである。

\begin{lemma}[有限次元 affine 加群の coherent な延長]
\label{cmo:lem:chart-ext}
$A$ を連結局所有限な次数付き右 noetherian 代数とし、
$u\in Z(A)_1$ を正則元とする。
$L=A[u^{-1}]$、$D=L_0$ と置く。
$D$ 上の有限次元右加群 $M$ に対し、
\[
 P(M)=\bigoplus_{n\ge0}M u^n,
 \qquad (m u^n)a_d=(m(a_du^{-d}))u^{n+d}
\]
を次数付き $A$-加群とし、$\mathcal P(M)=\pi P(M)$ と置く。
これは coherent であり、有限次元 $D$-加群 $M,N$ に対し全ての $i\ge0$ で
\[
 \operatorname{Ext}^i_{\operatorname{qgr}A}(\mathcal P(M),\mathcal P(N))
 \simeq\operatorname{Ext}^i_D(M,N)
\]
が成立する。
\end{lemma}
\begin{proof}
$u$ は中心で次数一の可逆元なので $L\simeq D[u,u^{-1}]$ であり、
次数零部分を取る関手は $\operatorname{Gr}L\simeq\operatorname{Mod}D$ を与える。
平坦な中心局所化 $\operatorname{Gr}A\to\operatorname{Gr}L$ は exact であり、
torsion 加群を消すので exact な関手
\[
 j^*:\operatorname{QGr}A\longrightarrow\operatorname{Mod}D
\]
を誘導する。

その右随伴は
$j_*N=\pi\operatorname{Res}(N\otimes_\C\C[u,u^{-1}])$ である。
ここで adjunction が Serre 商に降りることを確認する。
$V\in\operatorname{Gr}L$ の $L$ 上の単射分解を restriction すると、
局所化が exact であるため $\operatorname{Gr}A$ 上の単射分解になる。
従って任意の torsion $T$ に対し
\[
 \operatorname{Ext}^q_{\operatorname{Gr}A}(T,\operatorname{Res}V)
 \simeq\operatorname{Ext}^q_{\operatorname{Gr}L}(T[u^{-1}],V)=0
 \quad(q\ge0).
\]
ゆえに $\operatorname{Res}V$ は torsion-local であり、
商写像の section と同一視できる。これで上の adjunction を得る。
$j_*$ は $\pi$, restriction, $N\mapsto N[u,u^{-1}]$ の合成なので exact であり、
$j^*j_*\simeq\operatorname{id}$ である。
さらに exact な左随伴 $j^*$ を持つので $j_*$ は単射加群を単射対象へ移す。
$N$ の単射分解を $j_*$ で移して計算すれば
\[
 \operatorname{Ext}^i_{\operatorname{QGr}A}(j_*M,j_*N)
 \simeq\operatorname{Ext}^i_D(M,N)
\]
を得る。

$M$ が有限次元なら $P(M)$ は次数零の有限個の元で生成される。
$P(M)\subset M[u,u^{-1}]$ の商は負次数だけにあり、
各斉次元は十分高い正次数部分により消されるので torsion である。
従って $j_*M=\pi P(M)=\mathcal P(M)$ は coherent である。

最後に $\operatorname{QGr}A$ は局所 noetherian Grothendieck 圏で、
その noetherian 対象の部分圏が $\operatorname{qgr}A$ である。
この包含は全ての Ext を保つ。
これは Positselski--Schn\"urer \cite[Proposition~2.1]{CMOSmallBig} の適用である。
同命題の仮定は、noetherian 対象への任意の全射について、その始域の
有限生成部分対象を取ってもなお全射となることから従う。
従って $D^-(\operatorname{qgr}A)\to D^-(\operatorname{QGr}A)$ は完全忠実であり、
上の比較を $\operatorname{qgr}$ に制限できる。
\end{proof}

\subsection{有理係数の非滑らかな混合中心商}

\begin{proposition}[無限自己 Ext を持つ有理係数例]
\label{cmo:prop:rational-singular}
$\alpha=(9,9,1)$ とし、独立な二組の変数を持つ
$\mathcal B_\alpha$ を中心変数 $u$ で tensor する。
\[
 B=\mathcal B_\alpha\otimes_{\C[u]}\mathcal B_\alpha',\qquad
 A=B\big/(\Omega_\alpha+100\Omega_\alpha'-\tfrac{1616}{3}u^3)
\]
と置く。$p=p'=(3,-3,1)$ の定常 character に対応する
$\mathcal P\in\operatorname{qgr}A$ に対し
\[
 \sum_{j\ge0}\dim_\C\operatorname{Ext}^j_{\operatorname{qgr}A}
       (\mathcal P,\mathcal P)\,s^j
 =\frac{(1+s^3)^2}{1-s^2}.
\]
従って $\operatorname{qgr}A$ の homological dimension は無限で、
$D^b(\operatorname{qgr}A)$ は proper でない。
線束生成部分 $\mathcal C_A=\operatorname{thick}\{\mathcal O(j):j\in\mathbb Z\}$ は
proper であるが smooth でなく、線束生成残余部分 $\mathcal K_A^{\rm lin}$ も smooth でない。
また $\mathcal P\notin\mathcal C_A$ なので
$D^b(\operatorname{qgr}A)/\mathcal C_A\ne0$ である。
\end{proposition}
\begin{proof}
本節の普遍 PBW 表示と中心元の証明により、$B$ は左右 noetherian な
Koszul AS-regular 代数で、global dimension と parameter は $7$ である。
中心三次元は正則で、$u$ も $A$ 上正則である。
後者二点は $u=0$ での $\Omega_0+100\Omega_0'$ の正則性からも直接従う。
実際、各 $S=S(1,-1,-1)$ は正則中心元 $\Omega_0$ 上自由であり、
$S\otimes S$ は $\C[\Omega_0,\Omega_0']$ 上自由なので、この中心線型和は正則である。
$B$ の $u$-進 lowest term で中心元の正則性を持ち上げ、
$u x=w y$ を $u$ で還元して $y=u z$ とすれば、$u$ の商上の正則性も得る。

指定 character では
\[
 (ab)^3+(bc)^3+(ca)^3=-729,\qquad
 9a^2b^2c^2=729,\qquad
 \kappa(p)=0,\quad\lambda_p=\frac{16}{3}.
\]
従って二因子の中心和はこの character を消し、自己 Ext の計算複体は
$\C[t]\otimes\bigwedge(\eta_1,\eta_2)$、微分零となる。
前補題によって affine 自己 Ext は全次数で coherent 自己 Ext と一致する。
これが表示した Hilbert 級数と最初の二つの結論を与える。

補題\ref{crs:lem:sections}と\ref{crs:lem:finite-lines}の仮定は満たされ、
$\mathcal C_A$ は $\bigoplus_{j=0}^6\mathcal O(j)$ で生成され proper である。
もし smooth なら定理\ref{crs:thm:coherent}により
$\mathcal C_A=D^b(\operatorname{qgr}A)$ となり、右辺も proper になるが、
$\mathcal P$ の無限自己 Ext と矛盾する。
さらに命題\ref{crs:prop:residual-smooth}の対偶から残余線束部分も smooth でない。
$\mathcal P$ 自身が proper な $\mathcal C_A$ に入ることも同じ理由で不可能である。
\end{proof}

\paragraph{結論の範囲.}
この例は滑らかな非可換三次多様体の仮定を満たさないため、
滑らかな対象についての可換 Ku 実現予想への反例ではない。
また $D^b(\operatorname{qgr}A)$ の非proper性と、同 dg 圏の smoothness 単独は区別する。
従って本命題は、全 coherent 圏そのものの dg smoothness 単独を否定する主張へ
言い換えない。


\subsection{定数型 Sklyanin 変形の普遍中心と有限中心性}
\label{cmo:subsec:pi-center}

本節では斉次 Sklyanin 入力を $S=S(1,-1,-1)$ に固定する。
その定数項を任意に変えた PBW 変形について、中心の生成元を明示する。
以下で使う $a,b,c$ は一ブロックの三つの変形係数であり、
中心和の切断係数とは区別する。

\begin{lemma}[普遍六次中心元]\label{cmo:lem:universal-sixth-center}
中心変数 $a,b,c$ を用いて
\[
 \mathcal D=\C[a,b,c]\langle x,y,z\rangle/
 \bigl([y,z]-x^2+a,\ [z,x]-y^2+b,\ [x,y]-z^2+c\bigr)
\]
と置く。$a,b,c$ に重み二、$x,y,z$ に重み一を与える。
次の元は $\mathcal D$ の中心に属する：
\begin{align}
 C_0={}&z(z^2-c)(xzy+ax+by)+abz^2,
 \label{cmo:eq:center-six-zero}\\
 C_+={}&yz^4y+zyz^3y+z^2y^2z^2+2z^2yz^2y-z^4y^2+z^5x
 \notag\\
 &+a(zyz^2+2z^2yz+2z^3y)
   +b(-yz^2x-2zy^3+4z^2yx+4z^4)
 \notag\\
 &+c(-y^2z^2-3yz^2y-z^2y^2-2z^3x)
 \notag\\
 &+a^2z^2+b^2(-yz+3zy)+c^2(2y^2+zx)
 \notag\\
 &+2abzx+ac(-yz-2zy)+bc(-yx-5z^2).
 \label{cmo:eq:center-six-positive}
\end{align}
さらに、積の順序を逆転し、$x,a$ を固定し、
$y,z$ および $b,c$ をそれぞれ交換する反自己同型を $\iota$ と置き、
$C_-=\iota(C_+)$ とする。これも中心元である。

$a=b=c=0$ での像を $C_{0,0},C_{+,0},C_{-,0}\in S_6$ と書く。
既述の三次中心元
\[
 \Omega_0=\frac{x^3+y^3+z^3}{3}-\frac23\operatorname{sym}(xyz)
\]
について、四元
\begin{equation}\label{cmo:eq:sixth-center-basis}
 \Omega_0^2,\quad C_{0,0},\quad C_{+,0},\quad C_{-,0}
\end{equation}
は $S_6$ 内で一次独立である。
\end{lemma}

\begin{proof}
三つの関係を $r_x,r_y,r_z$ と書く。
$\iota(r_x)=r_x$、$\iota(r_y)=r_z$、$\iota(r_z)=r_y$ であるから、
$\iota$ は実際に $\mathcal D$ の反自己同型である。
従って $C_-$ の中心性は $C_+$ の中心性からも従う。

中心性と一次独立性には、次の有限有理数計算を用いる。
重み $d$ における自由代数の基底を
\[
 \bigl\{a^{e_1}b^{e_2}c^{e_3}w\ \bigm|\
       2(e_1+e_2+e_3)+|w|=d\bigr\}
\]
とし、全ての二側関係倍
$a^{e_1}b^{e_2}c^{e_3}v r_j w$ を行とする。
行簡約は $\mathbb Q$ 上で行うので、零への簡約は二側関係イデアルへの
厳密な所属証明である。重み六と七での行列の大きさ、rank はそれぞれ
\[
\begin{array}{c|r|r|r|r}
 d&\text{行数}&\text{列数}&\text{rank}&\text{商の次元}\\ \hline
 6&1476&1036&917&119\\
 7&5454&3108&2919&189
\end{array}
\]
となる。$[x,C_\nu],[y,C_\nu],[z,C_\nu]$ は
$\nu\in\{0,+,-\}$ の全てについて重み七の関係行列で零へ簡約される。
同じ計算は
$\Omega=\Omega_0-ax-by-cz$ の中心性も重み四で確認する。

active word の長さを降順、同じ長さでは $x<y<z$ の辞書式順序とし、
同じ word に付くパラメータ指数を辞書式に並べる。
重み六の簡約後に $a=b=c=0$ とすると、$S_6$ の二十八個の自由列を得る。
式\eqref{cmo:eq:sixth-center-basis}を行に取り、そのうち
\[
 y^2z^3y,\quad yz^2y^3,\quad yz^4y,\quad z^3y^3
\]
の四列を選んだ係数行列は
\begin{equation}\label{cmo:eq:sixth-center-minor}
 \begin{pmatrix}
  4/9&0&0&4/3\\
  0&0&0&-1\\
  0&0&1&0\\
  0&-1&0&0
 \end{pmatrix},\qquad \det=-\frac49\ne0
\end{equation}
である。よって一次独立性が従う。
この計算を再生成する有限証明書が
\path{computations/verify_mixed_pi_center.py} である。
係数の代入・補間・浮動小数点 rank による検査ではなく、
三つのパラメータを独立変数に保つ普遍恒等式を確認している。
\end{proof}

\begin{theorem}[全定数変形での有限中心性]\label{cmo:thm:finite-center}
任意の $\alpha=(\alpha_1,\alpha_2,\alpha_3)\in\C^3$ に対して
\[
 D_\alpha=\C\langle x,y,z\rangle/
 \bigl([y,z]-x^2+\alpha_1,\ [z,x]-y^2+\alpha_2,
       \ [x,y]-z^2+\alpha_3\bigr)
\]
と置く。$C_{\nu,\alpha}$ を上の $C_\nu$ の係数特殊化とすると
\begin{equation}\label{cmo:eq:affine-center-generators}
 Z(D_\alpha)=\C[\Omega_\alpha,C_{0,\alpha},C_{+,\alpha},C_{-,\alpha}].
\end{equation}
ここで affine 表示では
$\Omega_\alpha=\Omega_0-\alpha_1x-\alpha_2y-\alpha_3z$ である。
$D_\alpha$ はこの可換有限型中心環上有限であり、従って PI である。

さらに、$u$ を次数一の中心変数とする斉次化 $\mathcal B_\alpha$ で
\[
 \widehat C_{\nu,\alpha}
     =C_\nu(x,y,z;\alpha_1u^2,\alpha_2u^2,\alpha_3u^2)
\]
と定めると
\begin{equation}\label{cmo:eq:homogeneous-center-generators}
 Z(\mathcal B_\alpha)=\C[u,\Omega_\alpha,
             \widehat C_{0,\alpha},\widehat C_{+,\alpha},\widehat C_{-,\alpha}].
\end{equation}
右辺の $\Omega_\alpha$ は $\mathcal B_\alpha$ の斉次三次中心元を表す。
$\mathcal B_\alpha$ もこの中心上有限である。
\end{theorem}

\begin{proof}
まず、全次数に必要な既知の中心生成定理を特定する。
$S=S(1,-1,-1)$ の点曲線は
$E=V(x^3+y^3+z^3+xyz)$ であり、その点 shift は位数六である。
後者は Walton \cite[Proposition~4.1]{CMOWaltonRepresentations} にも
同じパラメータで記載されている。
Smith--Tate \cite[Theorem~3.7, $r=1$]{GRSklyaninST} によれば、
位数 $n$ の三次元 Sklyanin 代数の中心は
次数三の元 $g$ と次数 $n$ の三元 $z_1,z_2,z_3$ で生成される。
この定理は $3\mid n$ の場合も含む。
本件の $n=6$ に適用すると
\[
 Z(S)=\C[g,z_1,z_2,z_3],\qquad
 \deg g=3,\quad\deg z_i=6.
\]
従って $\dim_\C Z(S)_3=1$ かつ $\dim_\C Z(S)_6\le4$ である。
$\Omega_0$ は非零である：例えば可換化
$S_{\mathrm{ab}}=\C[x,y,z]/(x^2,y^2,z^2)$ への像は
$-2xyz/3\ne0$ である。よって $\Omega_0$ は $g$ の非零 scalar 倍である。
補題\ref{cmo:lem:universal-sixth-center}は $Z(S)_6$ の四つの独立元を与えるので、
\emph{$\dim_\C Z(S)_6=4$ であり、式\eqref{cmo:eq:sixth-center-basis}はその全体の基底}
となる。特に各 $z_i$ はこの四元の線形結合であり、
\begin{equation}\label{cmo:eq:graded-center-generators}
 Z(S)=\C[\Omega_0,C_{0,0},C_{+,0},C_{-,0}].
\end{equation}
これは低次数の中心計算だけから全次数の生成性を推測した結論ではない。
全次数の生成上界が Smith--Tate の定理であり、有限計算はその全生成元の
持上げを確認している。

既述の PBW 判定により $D_\alpha$ の active 濾過の随伴代数は $S$ である。
任意の中心元の最高部分は $Z(S)$ に属する。
式\eqref{cmo:eq:graded-center-generators}の全生成元は、補題で与えた
$D_\alpha$ の中心元へ持ち上がっている。
従って任意の $h\in Z(D_\alpha)$ の最高部分と一致する中心生成元の多項式を
$h$ から引き、active 次数について帰納すれば
式\eqref{cmo:eq:affine-center-generators}を得る。

同じ Smith--Tate の定理により $S$ は $Z(S)$ 上有限である。
有限個の斉次加群生成元を取り、これらを $D_\alpha$ に持ち上げる。
任意の元の最高部分をそれらで表し、次数を下げる帰納法を繰り返すと、
同じ有限個の持上げが $D_\alpha$ を $Z(D_\alpha)$ 上生成する。

斉次化については、$C_\nu$ の普遍恒等式に
$(a,b,c)=(\alpha_1u^2,\alpha_2u^2,\alpha_3u^2)$ を代入すれば
$\widehat C_{\nu,\alpha}$ の中心性を得る。
$u$ は $\mathcal B_\alpha$ 上正則であり、$\mathcal B_\alpha/(u)=S$ である。
これらの中心元の像が $Z(S)$ を全て生成するため、
Smith--Tate \cite[Lemma~3.6]{GRSklyaninST} が適用でき、
式\eqref{cmo:eq:homogeneous-center-generators}と加群有限性を得る。
実際、この段階の帰納法は、中心元 $h$ の像を生成元の多項式で消して
$h-P=u h_1$ と書き、$u$ の正則性から $h_1$ の中心性を得るものである。
\end{proof}

\begin{corollary}[係数族に一様な有限中心模型]
\label{cmo:cor:relative-finite-center}
上の中心元は三つの定数係数に関して多項式であり、係数の分母は生じない。
したがって、係数環 $R$ 上で $\alpha_i\in R$ を指定し、既述の分裂した
PBW 模型を底変更して得る $\mathcal B_{\alpha,R}$ は、
\[
 R[u,\Omega_\alpha,
       \widehat C_{0,\alpha},\widehat C_{+,\alpha},\widehat C_{-,\alpha}]
\]
という可換有限型中心部分環上有限である。
特に $R=\C[t]$、$\alpha\mapsto t^2\alpha$ の収縮族に適用できる。
複数因子が同じ $u$ を共有する tensor 積も、各因子のこの有限中心部分環の
$R[u]$ 上の tensor 積の像に対して有限である。
\end{corollary}

\begin{proof}
$S$ の $Z(S)$ 上の有限斉次加群生成元を一度選び、その同じ語を全係数族で使う。
PBW によって $\mathcal B_{\alpha,R}/(u)=S\otimes_\C R$ であり、
式\eqref{cmo:eq:graded-center-generators}の係数拡大は、表示した中心部分環の像に
含まれる。任意の斉次元を $u=0$ でこの有限個の語の結合として表し、差を $u$ で割る。
内部次数が毎回一ずつ下がるため、有限回で停止する。
この証明は $R$ の点ごとの有限生成性からの推論ではなく、一組の有限生成語を
全係数族で共通に用いる。
因子ごとの有限生成系の tensor 積を取れば最後の主張が従う。
\end{proof}

\begin{remark}[先行研究との関係と新規性の範囲]
\label{cmo:rem:pi-center-literature}
Walton の学位論文 \cite[Definition~V.1, Remark~V.16]{CMOWaltonThesis} は、
三次元 Sklyanin 代数の一般中心拡大を扱い、位数六の場合について
三次・六次中心元を持ち上げる方法を提案している。
本節はそのうち、斉次パラメータが $(1,-1,-1)$、線形補正 $d_i x_i u$ が零で、
定数補正 $e_i u^2$ が任意の部分族について、その持上げを明示し検証したものである。
位数六の全 Sklyanin パラメータや、非零の線形補正を含む一般拡大を扱う主張ではない。

Etingof--Ginzburg \cite[Theorem~3.3.2 and Section~9]{NFMEG} の
三次中心元と PBW 変形の理論、および Smith--Tate の斉次中心生成定理は既知の入力である。
本節の具体的六次持上げとこの定数型部分族への適用について、
後続文献を含む初出の網羅的確認はまだ完了していない。
有限中心性は明示した中心元から証明しており、associated graded の PI 性だけから
PBW 変形の PI 性を結論していない。
\end{remark}

\subsection{分離表示への代数的収縮と有限線束模型}

本節の混合表示には、中心係数の無限遠を分離表示へ戻す代数的収縮がある。
この収縮は実際の代数の同型を与えるが、それだけで
smoothness や可換 Ku 実現の開性を結論しない。

\begin{proposition}[混合係数の収縮と中心係数の三乗被覆]
\label{cmo:prop:contraction}
$k\geq2$ とし、$\alpha^{(i)}\in\C^3$、$a_i\in\C^\times$ を固定する。
例\ref{cmo:ex:sixty-four}の表示と同じ
$\mathcal B_{\alpha^{(i)}}$ を用い、
定理\ref{cmo:thm:finite-center}の普遍中心元
$\Omega_{\alpha^{(i)}}$ を用いる。
$R=\C[t]$、$\deg t=0$ と置き、
\[
 \mathscr B=
   \bigotimes_{R[u],\,i=1}^{k}\mathcal B_{t^2\alpha^{(i)}},
 \qquad
 w_t=\sum_i a_i\Omega_{t^2\alpha^{(i)}}+u^3,
 \qquad \mathscr A=\mathscr B/(w_t)
\]
とする。ここでは $x_i,y_i,z_i,u$ は内部次数一である。
各 $\mathcal B_{t^2\alpha^{(i)}}$ の式は
\[
 [y_i,z_i]=x_i^2-t^2\alpha^{(i)}_1u^2
\]
とその巡回置換、および $u$ の中央性で与える。

この族は以下を満たす。
\begin{enumerate}[label=(\arabic*)]
\item $\mathscr B$ は $R[u]$ 上自由であり、$w_t$ は正則中心三次である。
      $\mathscr A$ の各内部次数は有限自由な $R$-加群であり、
      各ファイバーの Hilbert 級数は
      \[
        \frac{1-s^3}{(1-s)^{3k+1}}
      \]
      である。
\item $S=\C\langle x,y,z\rangle/
      ([y,z]-x^2,[z,x]-y^2,[x,y]-z^2)$ と書くと、
      $t=0$ は分離表示
      \[
       \mathscr A_0=
       \bigl(S^{\otimes k}\otimes\C[u]\bigr)/
          \left(\sum_i a_i\Omega_0+u^3\right)
      \]
      である。従って定理\ref{csc:thm:comparison}の下で、
      これは滑らかな可換三次 $(3k-1)$-fold の全 coherent Ku 成分を持つ。
\item $t_0\ne0$ では、元の中心変数を $v$ と書くことにより
      \begin{equation}\label{cmo:eq:contraction-isomorphism}
       \mathscr A_{t_0}\simeq
       \left(\bigotimes_{\C[v],\,i=1}^{k}
                           \mathcal B_{\alpha^{(i)}}\right)
       \bigg/\left(\sum_i a_i\Omega_{\alpha^{(i)}}+t_0^{-3}v^3\right).
      \end{equation}
      同型は全 active 変数を保ち、$v\mapsto t_0u$ とする次数付き代数の同型である。
\end{enumerate}
\end{proposition}

\begin{proof}
元の混合周囲代数を $\mathcal B^{\mathrm{orig}}$ とし、
その中心変数を $v$ と書く。中心環の準同型
\[
 \C[v]\longrightarrow\C[t,u],\qquad v\longmapsto tu
\]
による底変更が、表示した $\mathscr B$ そのものである。
例\ref{cmo:ex:sixty-four}で使った Braverman--Gaitsgory の判定は、
三次重なりの定数補正について $\beta\otimes1-1\otimes\beta=0$ を使うだけなので、
任意の固定した $\alpha^{(i)}$ に対して PBW 性を与える。
その Rees 代数は $\C[v]$ 上自由であり、$k$ 因子の tensor 積も同様である。
従って $\mathscr B$ は $R[u]$ 上自由である。
各 $\Omega_{t^2\alpha^{(i)}}$ は指定した中心元の底変更なので中心であり、
$u$ も中心である。

$u=0$ では $\mathscr B/(u)=S^{\otimes k}\otimes R$ で、
$w_t$ は $\sum_i a_i\Omega_0$ となる。
各 $S$ は $\C[\Omega_0]$ 上自由なので、この和は正則である。
従って $w_t f=0$ なら、まず $u=0$ へ移して $f\in u\mathscr B$ を得る。
$u$ の正則性により同じ議論を繰り返せる。
固定内部次数では $u$ で無限回割れないので $f=0$ である。
同じ証明は任意の $t=t_0$ のファイバーで成立する。

各内部次数における $w_t$ の乗法は有限自由 $R$-加群間の写像であり、
全ての閉ファイバーで単射である。
最大小行列式による階数判定から、その余核は局所自由である。
$R=\C[t]$ は PID なので余核は自由であり、
正則中心三次の二項完全列から Hilbert 級数を得る。

$t=0$ の式は定義から従う。
$S$ の点曲線は $E=V(x^3+y^3+z^3+xyz)$、
$\Omega_0$ はその THCR 商を与える標準中心三次の非零倍である。
分離表示の中心式は $u^3$ で、その Jacobian 環は有限次元であるため、
定理\ref{csc:thm:comparison}が適用できる。

最後に $t_0\ne0$ なら、$v=t_0u$ によって全二次関係と
$\Omega_{\alpha^{(i)}}$ がそれぞれ
$\mathcal B_{t_0^2\alpha^{(i)}}$ と $\Omega_{t_0^2\alpha^{(i)}}$ の式へ移る。
また $t_0^{-3}v^3=u^3$ である。
逆変換 $u=t_0^{-1}v$ も次数を保つため、
式\eqref{cmo:eq:contraction-isomorphism}は実際の代数同型となる。
\end{proof}

\begin{proposition}[収縮族の proper な有限線束模型]
\label{cmo:prop:proper-contraction}
上の族で $N=3k+1$、$n=N-2$、$a=N-3$ と置く。
$\mathscr B$ は左右 noetherian で、$R$ に相対的な長さ $N$ の
Koszul 双加群解消とその片側 augmentation 解消を持つ。
それらは任意の係数変更と両立し、相対 Gorenstein 型は
\[
 R\underline{\operatorname{Hom}}_{\mathscr B}(R,\mathscr B)
       \simeq R(N)[-N],\qquad
 R\underline{\operatorname{Hom}}_{\mathscr A}(R,\mathscr A)
       \simeq R(a)[-(N-1)]
\]
である。

$\mathcal O(j)=\pi\mathscr A(j)$ とすると、有限線束部分
$\mathscr C^{\mathrm{line}}$ とその残余部分は
\[
 E=\bigoplus_{j=0}^{N-1}\mathcal O(j),\qquad
 \mathscr C^{\mathrm{line}}=\operatorname{thick}_R(E),\qquad
 G=L_{\mathcal O}\cdots L_{\mathcal O(a-1)}(E)
\]
で生成される。
特に
\[
 \Gamma_R=R\operatorname{End}(G)\in\operatorname{Perf}(R),\qquad
 \mathscr K^{\mathrm{line}}\simeq\operatorname{Perf}(\Gamma_R)
\]
という proper な相対 dg 模型を持つ。
 $R$ の局所化および閉ファイバーへの通常の底変更は、
この線束模型と残余模型を対応する代数の模型へ移す。
相対 smoothness は、定理\ref{cmo:thm:relative-smoothness}で有限中心模型から導く。
\end{proposition}

\begin{proof}
各因子は $Q=R[u]$ 上の cyclic potential
\[
 \Phi_i=x_iy_iz_i-y_ix_iz_i
       -\frac{x_i^3+y_i^3+z_i^3}{3}
       +t^2u^2\left(\alpha^{(i)}_1x_i+
                     \alpha^{(i)}_2y_i+\alpha^{(i)}_3z_i\right)
\]
の三つの cyclic 微分で定まる。
active 濾過の随伴代数は $S\otimes Q$ であり、PBW 分裂は上の底変更と両立する。
定理\ref{thm:nfm-relative-ambient-cy}の証明中の、
中心環に相対的な長さ三の potential 解消を各因子で作る。
その $Q$ 上の tensor 積は長さ $3k$ の相対解消であり、
完全性は分裂した PBW 濾過の随伴複体で確認できる。
$C=\mathscr B\otimes_R\mathscr B^{\mathrm{op}}$、
$D=\mathscr B\otimes_Q\mathscr B^{\mathrm{op}}$ と置く。
共有する中心変数 $u$ の左右差 $z=u_L-u_R$ は $C$ 上正則で、$D=C/(z)$ である。
その Koszul 双対は
$R\underline{\operatorname{Hom}}_C(D,C)=D(1)[-1]$ を与える。
従って potential 解消の自己双対性と change of rings により
\[
 R\underline{\operatorname{Hom}}_C(\mathscr B,C)
 \simeq R\underline{\operatorname{Hom}}_D
              (\mathscr B,D(1)[-1])
 \simeq\mathscr B(N)[-N]
\]
を得る。ここでは中心差の因子を一度だけ使っており、
異なるブロックで共有する $u$ を重複して数えていない。
また $D$ は $C$-perfect、$\mathscr B$ は $D$-perfect なので、
$\mathscr B$ は $C$-perfect である。

二次 Koszul 模型はこの change of rings とは別に指定できる。
各体値で active 濾過の随伴代数は
$S^{\otimes k}\otimes\C[u]$ であり、有限な内部次数の bar spectral sequence により
その Koszul 性が元の代数へ移る。
Hilbert 級数 $(1-s)^{-N}$ と合わせると、Koszul 空間の階数は
$(1+s)^{3k}(1+s)=(1+s)^N$ の係数である。
被約基底 $R$ 上の一定階数判定によりこれらは有限射影加群で、
標準 Koszul 双加群複体は係数変更と両立して完全となる。
augmentation を施せば最初の Gorenstein 式となり、
正則中心三次の二項解消から第二式を得る。

$S^{\otimes k}$ の noetherian 性は補題\ref{mbnp:lem:noetherian}による。
$\mathscr B/(u)=S^{\otimes k}\otimes R$ は中心多項式拡大なので noetherian である。
正次数中心元 $u$ による初項の持上げは、固定内部次数では有限回で止まるので、
$\mathscr B$ も左右 noetherian となる。$\mathscr A$ もその商として同様である。

$R=\C[t]$ の大域次元は一であり、上の有限双加群解消を片側加群へ tensor すると
$\operatorname{gldim}\mathscr B\leq N+1$ を得る。
正則中心元の change of rings により、$\mathscr A$ の左右 injective dimension も有限である。
従って noetherian 性と相対 Gorenstein 式を合わせると、
Burke--Stevenson の $\chi$ と有限 cohomological dimension の入力も満たされる。
以上の有限解消と Gorenstein 型を用いると、
補題\ref{lp:lem:local-cohomology}と式\eqref{lp:eq:H}の証明が適用できる。
実際、ここで使う入力は上の noetherian 性・有限次元性・相対 augmentation 双対と次数別自由性であり、
特殊ファイバーの可換性ではない。
従って
\[
 \operatorname{Ext}^p(\mathcal O(i),\mathcal O(j))=
 \begin{cases}
  \mathscr A_{j-i},&p=0,\\
  (\mathscr A_{i-j-a})^\vee,&p=n,\\
  0,&\text{その他}
 \end{cases}
\]
となる。右辺の各次数は有限自由 $R$-加群である。
$R$ の ambient augmentation Koszul 解消を $\mathscr A$ へ導来 tensor する。
中心三次の二項解消から $\operatorname{Tor}^{\mathscr B}_i(R,\mathscr A)$ は
$i=0,1$ にのみあり、それぞれ $R,R(-3)$ である。
これらは内部次数について torsion なので、得られる有限複体は tails 内で完全となる。
両端は一つの線束で、中間項の shift は $1,\ldots,N-1$ にあるため、
任意の連続した $N$ 個の twist からその前後を生成できる。
これで $E$ の生成性と properness を得る。
命題\ref{lp:prop:residual}と同じ有限 mutation によって $G$ を作ると、
その自己 Hom は $R$-perfect である。
底変更についても命題\ref{lp:prop:basechange}の
次数別に安定する局所 cohomology と有限 Hom の議論を使う。
指定した局所化と閉ファイバーでは、$\mathscr B,\mathscr A$ の noetherian 性は
中央局所化または商から従う。底も有限大域次元の正則環であり、
有限解消と次数別自由性は保たれるので、その適用条件を満たす。
\end{proof}

\subsection{有限中心上の通常の加群に対する局所持上げ}

\begin{lemma}[中心的な一変数に沿う射影次元の持上げ]
\label{cmo:lem:projective-lift}
$(C,\mathfrak m)$ を noetherian な可換局所環、
$t\in\mathfrak m$ とする。$B$ は中央 $C$-代数で $C$ 上有限、
$M$ は有限 $B$-加群、$d\ge0$ とする。
$t$ は $B$ と $M$ 上正則であると仮定する。
このとき
\[
 \operatorname{pd}_{B/tB}(M/tM)\le d
 \quad\Longrightarrow\quad
 \operatorname{pd}_B M\le d.
\]
\end{lemma}
\begin{proof}
$B$ は左右 noetherian である。$M$ の有限階数自由分解の最初の
$d$ 段を取り、その第 $d$ syzygy を $K$ とする。
全ての syzygy は $t$-torsionfree であり、$t$ で割っても
分解の完全性が保たれる。従って $K/tK$ は $B/tB$ 上有限射影である。

$C$ の $(t)$-adic 完備化 $\widehat C$ は忠実平坦である
（\cite[Tag~00MC]{CMOStacksCompletion}）。
$\widehat B=B\otimes_C\widehat C$ と
$\widehat K=K\otimes_C\widehat C$ を取る。
有限性によりこれらは $(t)$-adically complete で、
$t$ は $\widehat B$ の Jacobson radical に属する。
$K/tK$ を表示する有限行列の idempotent は $\widehat B$ へ持ち上がる。
この非可換行列に対する事実は、$e^2-e\in(t)$ から出発し、
\[
 e\longmapsto e-(2e-1)^{-1}(e^2-e)
\]
で誤差を二乗化して極限を取れば直接分かる。
逆元は $(2e-1)^2=1+4(e^2-e)$ から存在する。
従って $K/tK$ を持ち上げる有限射影 $\widehat B$-加群 $Q$ がある。

$Q$ の射影性により、指定した $Q/tQ\simeq\widehat K/t\widehat K$ を
射 $Q\to\widehat K$ へ持ち上げる。
余核は $t$ で割ると零なので Nakayama によりこの射は全射である。
核 $L$ は有限であり、$\widehat K$ 上の $t$ の正則性から
$L/tL=0$ を得る。再び Nakayama により $L=0$、
従って $\widehat K$ は射影である。

有限射影性はこの中央の忠実平坦底変更から降下する。
具体的には、有限自由表示 $0\to J\to B^r\to K\to0$ の分裂障害を考える。
$K$ の有限階数自由分解と平坦性により
\[
 \operatorname{Ext}^1_B(K,J)\otimes_C\widehat C
 \simeq
 \operatorname{Ext}^1_{\widehat B}(\widehat K,\widehat J)=0.
\]
忠実平坦性から障害は零で、$K$ は有限射影である。
\end{proof}

\begin{lemma}[中央スキーム上の perfect locus]
\label{cmo:lem:perfect-locus}
$Y$ を noetherian スキーム、$\mathcal B$ を coherent な
$\mathcal O_Y$-代数、$M$ を coherent $\mathcal B$-加群とする。
$M_y$ が $\mathcal B_y$ 上有限射影次元を持つ点の集合は開である。
この開集合上では $M$ は局所的に有限生成射影分解を持つ。
\end{lemma}
\begin{proof}
アファイン近傍上で $M$ の有限自由分解を必要な段階まで取り、
射影になった syzygy を $K$ とする。
$K_y$ の有限自由加群からの射影表示とその splitting は有限個の元で記述される。
$K$ の有限表示性により、射と splitting の等式の分母を有限個払えば
$y$ の中央開近傍上で $K$ 自身が射影となる。
\end{proof}

\begin{proposition}[proper な中央台からの相対 perfect 性]
\label{cmo:prop:proper-perfect-spread}
$R=\mathbf C[t]$、$Y\to\operatorname{Spec}R$ を proper な
noetherian スキームとする。
$\mathcal B$ を coherent $\mathcal O_Y$-代数、$M$ を coherent
$\mathcal B$-加群とし、両者が $R$-flat であると仮定する。
$M_0$ が $\mathcal B_0$ 上 perfect なら、
$0$ を含む Zariski 開近傍 $U\subset\operatorname{Spec}R$ が存在し、
$M_U$ は $\mathcal B_U$ 上 perfect となる。
\end{proposition}
\begin{proof}
$Y_0$ の各点で前二補題を適用する。
非 perfect locus は閉集合で $Y_0$ と交わらない。
proper な射によるその像も閉であり $0$ を含まないから、
この像を除けばよい。
\end{proof}


\subsection{有限中心模型による収縮族の相対 smoothness}

\begin{theorem}[分離ファイバーを含む相対 smooth proper 近傍]
\label{cmo:thm:relative-smoothness}
命題\ref{cmo:prop:contraction}の任意の
$\alpha^{(i)}\in\C^3$、$a_i\in\C^\times$、$k\ge2$ を取る。
$0$ を含む Zariski 開近傍 $U\subset\mathbb A^1_t$ 上で、
命題\ref{cmo:prop:proper-contraction}の線束模型と残余模型は
相対 smooth proper である。特に
\begin{equation}\label{cmo:eq:relative-diagonal}
 \Gamma_U\in\operatorname{Perf}\!\left(
       \Gamma_U^{\mathrm{op}}\otimes_U^{\mathbf L}\Gamma_U\right).
\end{equation}
ここで追加の対角 perfect 性の仮定は課さない。
\end{theorem}
\begin{proof}
系\ref{cmo:cor:relative-finite-center}によって、$\mathscr A$ は
$u$、各 $\Omega_{t^2\alpha^{(i)}}$、各六次中心元の像が生成する
可換有限型次数付き中心部分環 $C_R$ 上有限である。
その生成元の内部次数は $1,3,6$ である。
$C_R\subset\mathscr A$ は $R=\C[t]$ 上 torsionfree であり、
$R$ が PID なので $R$-flat である。
六次 Veronese $C_R^{(6)}$ は次数一で生成される。
実際、重み $1,3,6$ の単項式で重みが六の倍数なら、重み六の単項式へ分割できる。

\[
 Z=\operatorname{Proj}C_R^{(6)},\qquad
 \mathcal A=\widetilde{\mathscr A^{(6)}}
\]
と置く。$Z\to\operatorname{Spec}R$ は projective かつ flat であり、
$\mathcal A$ は coherent な $\mathcal O_Z$-代数で $R$-flat である。
flatness は、中央元を反転した次数零部分が $R$-torsionfree であることから従う。
Veronese の tails 同値と、可換標準次数環上の Serre 同値を
有限 $C_R^{(6)}$-代数の加群作用とともに用いれば
\[
 \operatorname{QGr}\mathscr A\simeq\operatorname{Qcoh}(Z,\mathcal A),\qquad
 \operatorname{qgr}\mathscr A\simeq\operatorname{coh}(Z,\mathcal A)
\]
を得る。ここでは $\mathscr A$ が内部次数一で生成されることを使う。
$C_R$ のファイバーが全中心と一致するという仮定は不要である。

全ての元の twist $\mathcal O(j)$ が $\mathcal A$-perfect であることを確認する。
$Z$ を六次の中央元 $f$ の標準開集合で覆い、
$L=\mathscr A[f^{-1}]$ とする。
$f$ を長さ六の語の和として表示し、先頭または末尾の一文字を分けると
\[
 1\in L_1L_{-1},\qquad 1\in L_{-1}L_1
\]
となる。これを反復すれば $L$ は強 $\mathbb Z$-次数付きである。
従って $L_j$ は $L_0$ 上有限生成射影で、各 twist は局所射影である。
中央の ample twist $\mathcal A\otimes\mathcal O_Z(m)$ は
元の $\mathcal O(6m)$ に対応する。
これらは $D_{\mathrm{qc}}(Z,\mathcal A)$ を compact に生成する。
具体的には、忘却関手との随伴と ample twist の検出性
\cite[Tag~0BQR]{CMOStacksDerived}を用いる。
生成性は ample な切断によるアファイン局所化、compact 性は有限アファイン被覆の
\v Cech 複体と局所射影性から従う。
命題\ref{cmo:prop:proper-contraction}の有限線束和 $E$ は全 twist を厚生成するので、
$E$ はこの非有界導来圏の compact generator である。
この記述は閉ファイバーと $R$ の局所化に両立する。

$Y=Z\times_R Z$ 上で
\[
 \Lambda=\mathcal A^{\mathrm{op}}\boxtimes_R\mathcal A,
 \qquad M=\Delta_*\mathcal A,
 \qquad F=E^\vee\boxtimes_R E,
 \qquad E^\vee=R\mathcal Hom_{\mathcal A}(E,\mathcal A)
\]
と置く。$\Lambda,M$ は coherent かつ $R$-flat であり、
$M_0=\Delta_{0*}\mathcal A_0$ である。
双対は左右の perfect 圏の反同値なので、$E^\vee$ は反対側を生成する。
$F$ の厚生成部分は全ての中央 twist $\Lambda(i,j)$ を含み、
特に ample な $\Lambda(m,m)$ を含む。
従って $F$ は $D_{\mathrm{qc}}(Y,\Lambda)$ の compact generator である。

$\Sigma_R=R\operatorname{End}_{\mathcal A}(E)$ と書く。
局所有限射影加群の Hom は $R$-flat であり、有限アファイン被覆の
\v Cech 模型を tensor すると K\"unneth 同型を得る
\cite[Tag~0FU4]{CMOStacksDerived}。
局所的な有限射影加群の評価同型と合わせて
\begin{equation}\label{cmo:eq:kernel-diagonal}
 \begin{split}
 R\operatorname{End}_\Lambda(F)
      &\simeq\Sigma_R^{\mathrm{op}}\otimes_R^{\mathbf L}\Sigma_R,\\
 R\operatorname{Hom}_\Lambda(F,M)&\simeq\Sigma_R
 \end{split}
\end{equation}
となる。第二式は局所的に $T\mapsto[\varphi\otimes e\mapsto\varphi(T(e))]$ で与えられる。
$S^\vee\boxtimes U$ による前合成は $T\mapsto STU$ に対応するため、
左右の作用を保持し、$\Sigma_R$ の対角双加群そのものを与える。
compact generator による dg Morita 同値
\cite[Theorem~3.8(c)]{NFMKellerDg}から、
$M$ の $\Lambda$-perfect 性は線束圏の相対 smoothness と同値である。
同じ説明が $t=0$ でも成立する。

特殊ファイバーは分離表示であり、定理\ref{csc:thm:comparison}によって
全 coherent 圏が smooth proper である。
その有限線束和は生成対象なので、式\eqref{cmo:eq:kernel-diagonal}の特殊化により
$M_0$ は $\Lambda_0$-perfect である。
$Y\to\operatorname{Spec}R$ は proper なので、
命題\ref{cmo:prop:proper-perfect-spread}を $\Lambda,M$ に適用できる。
従って $0$ を含む開近傍 $U$ で $M_U$ は perfect であり、
線束圏は $U$ 上 smooth となる。
proper 性は既に命題\ref{cmo:prop:proper-contraction}で確認した。
最後に、有限例外列による admissible な残余成分には smooth proper 性が移るため、
$\Gamma_U$ について式\eqref{cmo:eq:relative-diagonal}を得る。
\end{proof}

\begin{corollary}[定数型混合族の余有限可換実現]
\label{cmo:cor:cofinite}
任意の $k\ge2$、$\alpha^{(i)}\in\C^3$、$a_i\in\C^\times$ を固定し、
命題\ref{cmo:prop:contraction}の定数型混合族を考える。
このとき元の $c$-族について有限集合 $F\subset\C$ が存在し、
$c\notin F$ では $D^b(\operatorname{qgr}A_c)$ は smooth proper であり、
ある滑らかな可換三次 $(3k-1)$-fold $X_c$ と
\[
 \operatorname{Ku}_{\mathrm{coh}}(A_c)
             \simeq\operatorname{Ku}(X_c)
\]
が存在する。
これは通常の dg Morita 同値であり、rotation・周期・point scheme の
標識の保存を追加しない。
\end{corollary}

\begin{proof}
定理\ref{cmo:thm:relative-smoothness}によって $\operatorname{Perf}(\Gamma_U)$ は
実際の相対 smooth proper な圏族である。
$n=3k-1\geq5$ なので $n\ne4$ であり、
$t=0$ の可換実現に定理\ref{grl:thm:open-locus}を適用できる。
従って $0$ を含む開近傍 $U'\subset U$ 上で
各残余線束圏は可換 Ku 成分を持つ。
各閉ファイバーの周囲代数は、上の有限 Koszul 双加群解消により大域次元が
高々 $N$、augmentation 双対の非零な最高 Ext により丁度 $N$ であり、
Gorenstein parameter $N$ の両側 noetherian Koszul AS-regular 代数である。
従って中心三次商に対する命題\ref{crs:prop:residual-smooth}と
定理\ref{crs:thm:coherent}の仮定を満たす。
残余圏と有限例外列の proper な gluing から線束部分の smoothness を得て、
その後の coherent 飽和により線束部分と全 coherent 圏が一致する。

$\mathbb A^1\setminus U'$ は有限集合である。
$t\mapsto t^{-3}$ は $\mathbb G_m$ の有限全射であり、
式\eqref{cmo:eq:contraction-isomorphism}により
$U'\cap\mathbb G_m$ の像に属する全 $c$ で結論を得る。
例えば $F$ は $0$ と $\mathbb A^1\setminus U'$ の非零点の像を含む
有限集合として取れる。
\end{proof}

\begin{remark}[有限中心性の役割と残る個別判定]
\label{cmo:rem:cofinite-scope}
この証明は、明示的な六次中心元による有限中心性、通常の coherent 対角加群、
proper な中央台を使う。一般の proper dg 代数について、smooth な特殊ファイバーや
全 Artin 切断での互換同値だけから相対 smoothness を推論しているのではない。
その推論への反例と、通常の完備底で生成性が失われる機構は研究ノートに記す。

有限集合 $F$ はまだ有効に決定していない。
$F$ を定理\ref{cmo:thm:point-norm}の norm の根集合と同一視する理由はない。
norm は定常点の出現を検出するものであり、smoothness の判別式とは別である。
例\ref{cmo:ex:sixty-four}の六十四個の値、$c=0$、および任意の指定された
個別の $c$ における全体の smoothness・Ku 比較は、この余有限系だけでは判定されない。
また得られた可換標的について、point scheme と Poisson 零スキームの
同時比較や、元のブロック分解の保存はまだ主張していない。
\end{remark}
